\section{Theoretical Guarantees}
\label{sec:identifiability_theory}

This section gives the guarantees behind steps 2, 4, 5 and 7 of
Algorithm~\ref{alg:iasa}, and the proofs are in Appendix~\ref{app:identifiability_proofs}.
Steps 1--3 construct \(\At\) before any reading is fitted, and the analysis treats
\(\At\) as fixed, which requires the following.

\begin{assumption}[Declared model]
\label{ass:model}
\(\mathbf y\) follows~\eqref{eq:general_model} with \(\mathbf c\ge0\), where \(\mathbf y\)
is the readings minus a declared baseline field, and \(A(\omega)\), \(Q\), the components,
\(\mathcal F_0\) and the set of observed rows are fixed independently of
\(\boldsymbol\epsilon\).
\end{assumption}

In the observed runs the baseline is the transported field of a map interpolated from
the first hour of readings, and its error enters \(\boldsymbol\epsilon\).  The assumption
fails if \(Q\) or the components are chosen after looking at estimates, or if the
transport model is wrong (Section~\ref{subsec:operator_error}).  Step 7 also requires a
noise model.

\begin{assumption}[Noise]
\label{ass:noise}
\(\boldsymbol\epsilon\sim N(0,\sigma^2I_N)\) (or, where stated, \(N(0,\Sigma)\)), and the
declared level satisfies \(\sigma_e\ge\sigma\) (or \(\sigma_e^2\ge\|\Sigma\|_2\)).
\end{assumption}

On the observed data the difference of the two co-located monitors has lag-1
autocorrelation \(0.857\), so independence in time does not hold, and
\(\sigma_e\ge\sigma\) is not checked (Section~\ref{subsec:observed_results}).

\subsection{Individual Coefficients}

We begin with the individual source--basis coefficients, such as the four time-of-day
activities of traffic.

\begin{proposition}[Exact identifiability]
\label{prop:rank_identifiability}
In the noiseless projected model \(\yt=\At\mathbf c\), \(\At\mathbf c_1=\At\mathbf c_2\)
implies \(\mathbf c_1=\mathbf c_2\) for all \(\mathbf c_1,\mathbf c_2\in\mathbb R_+^J\)
if and only if \(\operatorname{rank}(\At)=J\).
\end{proposition}

\begin{proposition}[Noise robustness]
\label{prop:noise_robust}
Assume \(\operatorname{rank}(\At)=J\) and let \(\sigma_J(\At)\) be the smallest
singular value.  If \(\widehat{\mathbf c}\) is the least-squares estimate, or the
nonnegative least-squares (NNLS) estimate and \(\mathbf c\ge0\), then
\begin{equation}
\|\widehat{\mathbf c}-\mathbf c\|_2\le\|P_{\At}\widetilde{\boldsymbol\epsilon}\|_2/\sigma_J(\At)
\le\|\widetilde{\boldsymbol\epsilon}\|_2/\sigma_J(\At).
\label{eq:nnls_noise_bound}
\end{equation}
\end{proposition}

\begin{remark}[Implication: coefficients]
Both propositions require only Assumption~\ref{ass:model}.  In every observed week
\(\operatorname{rank}(\At)=7=J\).  With noise, \(1/\sigma_J\) is the largest factor by
which noise enters the coefficients, so an identifiable problem can still have large
coefficient errors (Experiment 1).  IASA reports \(\sigma_J\) in step 4 and reports coefficients only for
single-component blocks.
\end{remark}

\subsection{Identifiable Groupings}

When the coefficients are not determined, the reported quantity is the signal of a
group of sources.  The following definition states when such signals are determined,
which is requirement R1.

\begin{definition}[Identifiable grouping]
\label{def:identifiable_grouping}
A partition \(\mathcal P\) of \(E\) is \emph{identifiable} if
\(\At\mathbf c_1=\At\mathbf c_2\) with \(\mathbf c_1,\mathbf c_2\in\mathbb R_+^J\)
implies \(\At_g\mathbf c_{1,g}=\At_g\mathbf c_{2,g}\) for every block
\(g\in\mathcal P\).
\end{definition}

\begin{theorem}[Identifiable groupings]
\label{thm:groupings}
\leavevmode
\begin{enumerate}[label=(\alph*),leftmargin=*,itemsep=0pt,topsep=2pt]
  \item \(\mathcal P\) is identifiable if and only if
  \(\operatorname{col}(\At)=\bigoplus_{g\in\mathcal P}\operatorname{col}(\At_g)\),
  equivalently \(\operatorname{rank}(\At)=\sum_{g\in\mathcal P}\operatorname{rank}(\At_g)\).
  \item There is a unique finest identifiable partition \(\mathcal P^\star\), and
  \(\mathcal P\) is identifiable if and only if every block of \(\mathcal P\) is a
  union of blocks of \(\mathcal P^\star\).
  \item Let \(B\subseteq E\) index a basis of \(\operatorname{col}(\At)\) chosen from
  the columns, and join each \(e\notin B\) to every \(b\in B\) whose coefficient is
  nonzero when \(\at_e\) is written in that basis.  The blocks of \(\mathcal P^\star\)
  are the connected components of this graph, for any choice of \(B\), and they are the
  connected components of the column matroid of \(\At\).
  \item For a declared partition \(\mathcal D\), the finest identifiable partition
  whose blocks are unions of blocks of \(\mathcal D\) is the finest common coarsening
  \(\mathcal P^\star\vee\mathcal D\).
\end{enumerate}
\end{theorem}

\begin{remark}[Implication: exact resolution]
The theorem requires only Assumption~\ref{ass:model}.  A partition meets R1 exactly when no
group's span meets the span of the other groups except at \(0\) (a), and there is one
finest such partition (b), computed before fitting in \(O(NJ^2)\) arithmetic (c).  A source
shares its block of \(\mathcal P^\star\vee\mathcal D\) with another source exactly when a
nonzero combination of its fingerprints equals a combination of the other
fingerprints.  In every observed week \(\mathcal P^\star\) is the seven single columns,
so the atoms are the four groups.  Step 5 forms its atoms by (c) and (d), so every
partition IASA can return meets R1.  For sums of coefficients, which are in different
proxy units, the minimal groupings can be non-unique and deciding whether one with at
least two blocks exists is NP-complete (Proposition~\ref{prop:sum_identifiable},
Appendix~\ref{app:identifiability_proofs}), so IASA reports signals instead.
\end{remark}

\subsection{Group-Signal Error}

Requirement R2 concerns the effect of noise, which identifiability does not control.
The error of a group's fitted signal is governed by its separation \(s_g\)
of~\eqref{eq:separation}, not by \(\sigma_J\).

\begin{theorem}[Group-signal error]
\label{thm:group_error}
Let \(\mathcal P\) be identifiable and let \(\widehat{\mathbf c}\) be any
least-squares solution, or any NNLS solution with \(\mathbf c\ge0\).
\begin{enumerate}[label=(\alph*),leftmargin=*,itemsep=0pt,topsep=2pt]
  \item The group signals \(\At_g\widehat{\mathbf c}_g\) do not depend on which
  solution is chosen, and for every \(g\in\mathcal P\),
  \begin{equation}
  \|\At_g(\widehat{\mathbf c}_g-\mathbf c_g)\|_2
  \le
  \frac{\|P_{\At}\widetilde{\boldsymbol\epsilon}\|_2}{s_g}.
  \label{eq:group_bound}
  \end{equation}
  \item If \(\boldsymbol\epsilon\sim N(0,\sigma^2I_N)\), for a least-squares solution and
  \(r_g=\operatorname{rank}(\At_g)\),
  \(\mathbb E\|\At_g(\widehat{\mathbf c}_g-\mathbf c_g)\|_2^2\le\sigma^2r_g/s_g^2\), with
  equality when \(r_g=1\).
  \item If \(\boldsymbol\epsilon\sim N(0,\sigma^2I_N)\), with probability at least
  \(1-\delta\), simultaneously for every identifiable partition and every block \(g\),
  \begin{equation}
  \|\At_g(\widehat{\mathbf c}_g-\mathbf c_g)\|_2
  \le
  \frac{\sigma\bigl(\sqrt{\operatorname{rank}(\At)}+\sqrt{2\log(1/\delta)}\bigr)}{s_g},
  \label{eq:group_bound_gauss}
  \end{equation}
  and if \(\boldsymbol\epsilon\sim N(0,\Sigma)\), the same holds with \(\sigma^2\)
  replaced by \(\|\Sigma\|_2\).
\end{enumerate}
\end{theorem}

\begin{remark}[Implication: the error bound IASA reports]
Part (a) holds for every noise vector under Assumption~\ref{ass:model}, whereas (b) and
(c) also need the Gaussian noise of Assumption~\ref{ass:noise}.  IASA, which uses NNLS, relies on
(a) and (c).  The columns inside a group, such as the four time-of-day columns of
traffic, can be nearly dependent, which makes the bound of
Proposition~\ref{prop:noise_robust} large while the group's signal error stays at most
\(\|P_{\At}\widetilde{\boldsymbol\epsilon}\|_2/s_g\).  With \(\sigma\) replaced by the
declared \(\sigma_e\), the right side of~\eqref{eq:group_bound_gauss} is the bound
\(b_g\) of step 7, which holds when \(\sigma_e\ge\sigma\) and the operator is exact.  Requiring \(b_g\le\eta\) for a tolerance \(\eta\) gives
the threshold \(\tau_\theta=\sigma_e(\sqrt{\operatorname{rank}(\At)}+\sqrt{2\log(1/\delta)})/\eta\)
of step 5, and \(0.141\) corresponds to \(\eta=7.1\,\sigma_e(\sqrt{\operatorname{rank}(\At)}+\sqrt{2\log(1/\delta)})\).
\end{remark}

\begin{definition}[Reportable partition]
\label{def:reportable}
A partition into unions of atoms is \emph{reportable at threshold
\(\tau_\theta\in(0,1]\)} if every block has \(s_g\ge\tau_\theta\).
\end{definition}

Reportable partitions are not closed under coarsening, and the minimal ones (with no
strictly finer reportable partition) need not be unique
(Proposition~\ref{prop:reportable_structure}), so with up to 12 atoms IASA checks every
union of atoms and lists ties.  The one-block partition is always reportable.

\subsection{Errors in the Operator}
\label{subsec:operator_error}

The operator is built from an estimated wind field and assumed dispersion parameters.
The next result bounds how far an error \(\Delta\) in \(\At\) moves the separations.

\begin{proposition}[Stability under operator error]
\label{prop:grouping_stability}
Let \(\At'=\At+\Delta\).  For \(S\subseteq E\) with \(\At_S\) of full column rank and
\(\|\Delta_S\|_2\le\sigma_{\min}(\At_S)/2\), let
\(\beta_S=\arcsin\bigl(\|\Delta_S\|_2/(\sigma_{\min}(\At_S)-\|\Delta_S\|_2)\bigr)\).
If this holds for \(S=g\) and \(S=E\setminus g\), with \(g\) nonempty and proper, then
\(|\theta_g(\At')-\theta_g(\At)|\le\beta_g+\beta_{E\setminus g}\).  If \(\At\) has full
column rank, \(\|\Delta\|_2\le\sigma_J(\At)/2\), and
\(|\theta_g(\At)-\arcsin\tau_\theta|>\beta_g+\beta_{E\setminus g}\) for every nonempty
proper union \(g\) of atoms, the reportable partitions of \(\At'\) and \(\At\) coincide,
and so do the minimal reportable partitions.
\end{proposition}

\begin{remark}[Implication: wind and dispersion error]
The proposition uses no noise assumption.  Under a wind or dispersion error \(\Delta\) that meets
the conditions, the set of reportable partitions is unchanged, and so is the reported
partition when one reportable partition has the most blocks.  In every observed week
each union of atoms has a separation angle of at least \(45.4^\circ\), against
\(\arcsin\tau_\theta=8.1^\circ\), but whether the true \(\Delta\) meets the conditions is not
known.
The fitted signals do change: for least squares, an operator error acts as additional
noise of size at most \(\|\Delta\|_2\|\mathbf c\|_2\)
(Proposition~\ref{prop:operator_perturbation}, Appendix~\ref{app:operator_error_details}),
which Experiment 5 measures.
\end{remark}
